%%%%%
%
%% Dans ce fichier, les �ventuels accents sont cod�s
%% en ISO-8859-15 (aka latin9)
%% La ligne suivante donne trois exemples 
%% e accent aigu : �, e accent grave : �, e-dans-l'o : �
\section{Correction de l'influence de la pression sur la mesure de Tdf}

\itemstyle{Op�rateur~:} \operfmt{Correc\_Tdf\_P}
\index{Correc\_Tdf\_P}

%%%%%OBJET%%%%
\itemstyle{Objet~:}
 Toute variation de pression d'une parcelle d'air change
proportionnellement la pression partielle de vapeur d'eau,
et donc la valeur du point de ros�e (ou de givre), qui
est la temp�rature telle que ewt(Td) = e.

Cet op�rateur corrige les mesures de Td/Tf faites � une pression
diff�rente de celle de l'air qu'on souhaite caract�riser.

%%%%%ARGUMENTS%%%%
\bigskip
%% Dans ce fichier, les �ventuels accents sont cod�s
%% en ISO-8859-15 (aka latin9)
%% La ligne suivante donne trois exemples 
%% e accent aigu : �, e accent grave : �, e-dans-l'o : �
%%%%%ENTREES%%%%
\begin{tabarg}
\multicolumn{3}{l}{\titretab {Entr\'ees}}\\
\hline
$Tdf_m$ & Vecteur, \texttt{float} &  Temp�rature de ros�e ou givre mesur�e (�C ou K)\\
$Ps_m$ & Vecteur, \texttt{float} & Pression de mesure (hPa)\\
$Ps$ & Vecteur, \texttt{float} & Pression de l'air (m�me unit� que $Ps_m$)\\
\end{tabarg}

\medskip
%%%%%SORTIES%%%%
\begin{tabarg}
\multicolumn{3}{l}{\titretab {Sorties}}\\
\hline
$Tdf$ & Vecteur, \texttt{float} & Temp�rature de ros�e ou givre corrig�e (m�me
unit� que $Tdf_m$)\\
\end{tabarg}


%%%%%FORMULE%%%%
\itemstyle{Formule}

%% Formule de l'operateur : Correc_Tdf_P
%% Dans ce fichier, les �ventuels accents sont cod�s
%% en ISO-8859-15 (aka latin9)
%% La ligne suivante donne trois exemples 
%% e accent aigu : �, e accent grave : �, e-dans-l'o : �
\[
Tdf = ew^{-1} (ew(Td_m) \frac{P}{P_m})
\]



%%%%%AUTEUR%%%%
\itemstyle{Auteur~:}
CNRM/GMEI/TRAMM

%%%%%REFERENCE%%%%
\itemstyle{R\'ef\'erences~:}
%INCONNU
